

\begin{lstlisting}[language=bash]
    brew install icarus-verilog
\end{lstlisting}

You may be able to generate bitstreams and load them onto boards; however, that's only half the FPGA design battle.
You also need to be able to simulate your designs. 
Luckily this is a piece of software which can be installed with homebrew; 
install Icarus Verilog with the command above. 


This will install iverilog to compile simulations which can then be run using vvp which was installed alongside iverlog.
We will go over the use of this momentarily; however, there is a final piece of the puzzle, viewing the waveforms.
You're going to want to export a vcd file from your testbenches.  

\begin{lstlisting}[language=verilog]
    initial begin
        $dumpfile("testbench.vcd");
        $dumpvars(0, testbench);
    end
\end{lstlisting}

Add these lines to your testbench. 
When the vvp simulator runs, it will generate a VCD file that you can view with a waveform viewer.


\begin{lstlisting}[language=bash]
    brew install surfer
\end{lstlisting}

The best option currently available is surfer which after running the above install can be launched from the command line with the surfer command. 
This will bring up the gui and allow you to import the vcd file and inspect the waveforms.